% Nouveau EGA V (cote 25): archive cote 25, site page 78 (PDF page 79).
% typescript. Draft compared: none.
% AI-transcribed by Claude Opus 5.5 (claude-opus-5-5, adaptive thinking, effort high), Message Batches API, batch ega25b (msgbatch_01A7n5YuXubUPNK1zsepi8RG), 2026-11-13; not yet checked by eye.
\scanpage{25/78}
\mspage{\hand{73}}
% NOTE: this typed line is framed by a handwritten curved line that leads down to the place after "$\underline{T}^{n-1}_Q$ ," in the next line, just before "soient"
\unsure{qui|qbi} sont aussi des sous-schémas fermés \inserted{disjoints} de $\underline{X}^{n-1}_Q$ ;

inverses $\underline{Y}^{n-1}_Q$ et $\underline{T}^{n-1}_Q$ , soient p,q,r les projections structurales
de ces schémas dans Q. Alors par \deleted{construction}\inserted{\unsure{hypothèse essentiellement}}, p est lisse en les
points au dessus de $\hand{\zeta}_o$ , et q est étale en les points au dessus de $\hand{\zeta}_o$ ;
il en est de même de r, car on voit aussitôt que $\underline{T}^{n-1}_Q$ n'est autre que
$Tx_kQ$ ( Q-isomorphisme). Enfin, le morphisme p est propre, et compte
tenu que X est géométriquement connexe, les fibres de p sont géométrique-%
ment connexes, (th de Bertini). Par suite, il existe un voisinage ouvert
V de $\hand{\zeta}_o$ dans X tel que \xout{\unsure{le morphisme}} $\underline{X}^{n-1}_Q|V = X'$ soit propre et lisse
sur V, à fibres géométriquement connexes, et comme la fibre de $\hand{\zeta}_o$ ,
qui n'est autre que $X\cap L_o$ , est de dimension 1, on peut supposer que les
fibres de X' sur V sont toutes de dim. 1. Enfin, prenant V assez petit,
on peut supposer que $\underline{Y}^{n-1}_Q|V$ et $\underline{T}^{n-1}_Q|V$ sont étales sur V, donc le
préschéma somme Z' de ses deux (qui s'identifie à un sous-\xout{\unsure{pré}}schéma
fermé de X' ) est étale sur V. Par suite, \xout{\illegible{1}} posant $U = X' - Z'$,\xout{\illegible{1}}
\xout{\illegible{1}} le morphisme $U\to V$ est un morphisme élémentaire. Mais \xout{\illegible{1}} U est aussi
un ouvert de \inserted{$X''=$}$\underline{X}^{n-1}_Q - \underline{Y}^{n-1}_Q - \underline{T}^{n-1}_Q$ , image inverse de X-Y-T dans $\underline{X}^{n-1}_Q$,
et d'autre part $X''\to$ X-Y-T est évidemment un isomorphisme (car
$\underline{X}^{n-1}_Q - \underline{T}^{n-1}_Q \to$ X-T est un isomorphisme ). Donc U s'identifie à un ouvert
de X-Y-T, ouvert contenant d'ailleurs $L_o\cap X$ et à fortiori x. C'est le
voisinage cherché de x contenu dans X-Y.

\mspage{\hand{78}}
